\begin{textsample}{POS Dim 2 – 2010\_c6 – Score 2.00 – 2010\_c6/enade\_2010\_ter\_ocup\_b707218cef.txt}  \label{ex:f2_pos_184}
CAPÍTULO 3 ANÁLISE TÉCNICA DA PROVA Este capítulo tem por objetivo apresentar o desempenho dos estudantes de Terapia Ocupacional no ENADE/2010.
Para isso, foram calculadas as estatísticas básicas da prova em seu todo, bem como as estatísticas dos componentes relacionadas à Formação Geral e ao de Conhecimento Específico da Área.
Nas tabelas, são apresentados o tamanho da população inscrita e de presentes, e as seguintes estatísticas das notas⁷ : média do desempenho na prova, erro padrão da média, desvio padrão, nota mínima, mediana, nota máxima, coeficiente de assimetria e de variação.
As estatísticas apresentadas neste capítulo contemplam o total de estudantes e, separadamente, os Ingressantes e os Concluintes.
Foram calculadas tendo -se em vista as seguintes agregações : ( a ) as Grandes Regiões e o país como um todo ; ( b ) a Categoria Administrativa ; e ( c ) a Organização Acadêmica.
Em relação aos gráficos de distribuição de notas, o intervalo considerado foi de 10 unidades, aberto à esquerda e fechado à direita, com exceção do primeiro intervalo, [ 0 ; 10 ], fechado em ambos os extremos.
Para os gráficos de distribuição das notas das questões discursivas, foram consideradas mais duas categorias : questão em branco e nota zero. 3.1 ESTATÍSTICAS BÁSICAS DA PROVA 3.1.1 Estatísticas Básicas Gerais A Tabela 3.1 apresenta as estatísticas básicas da prova por grupo de estudantes.
A população total de estudantes regularmente inscritos⁸ para a prova foi 2.054.
Destes, 1.808 estiveram presentes, sendo 12,0 \% o índice de não comparecimento.
A abstenção foi maior entre Ingressantes ( 17,3 \% ) do que entre Concluintes ( 4,0 \% ).
A nota média de toda a prova foi 40,1, sendo que os Ingressantes obtiveram média mais baixa, 37,7, que os Concluintes, cuja média foi 43,1.
O desvio padrão foi 13,7, sendo o do grupo de Ingressantes menor ( 12,4 ) do que o do grupo dos Concluintes ( 14,7 ) ; os coeficientes de variação ( 32,9 \% no grupo dos Ingressantes e 34,1 \% no dos Concluintes ) indicam maior dispersão relativa nas notas de Concluintes.
A nota máxima foi 74,9, obtida por Concluintes, ao passo que a maior nota obtida por Ingressantes foi 70,0. ⁷ Algumas dessas estatísticas estão definidas no Capítulo 1. ⁸ As estatísticas das notas apresentadas consideram apenas os estudantes regularmente inscritos e com resultados validados por ausência de erro administrativo.
Uma comparação entre o desempenho de Ingressantes e Concluintes também pode ser realizada a partir dos dados apresentados no Gráfico 3.1, com a distribuição de notas finais.
Destaca -se que as proporções mais elevadas de notas, considerando -se o intervalo posterior à nota 40,0, foram atingidas principalmente pelos Concluintes, cuja concentração das notas está à direita dos Ingressantes.
Já nas faixas de notas mais baixas, inverteu -se o quadro identificado nas faixas de notas mais altas.
Essa característica é corroborada pelo valor da nota mediana 38,8 no grupo dos Ingressantes, abaixo da encontrada no grupo dos Concluintes, 44,7, transparecendo uma superioridade de desempenho dos Concluintes em relação aos Ingressantes na prova em tela ( Tabela 3.1 ).
As faixas de notas modais foram ( 30 ; 40 ] para os Ingressantes com 29,6 \% dos estudantes deste grupo e ( 40 ; 50 ] para Concluintes com 27,9 \%.
Os coeficientes de assimetria indicam uma assimetria negativa para Ingressantes ( -0,4 ) e para os Concluintes ( -0,8 ).
O Gráfico 3.2 apresenta informações referentes à nota final de Ingressantes e Concluintes, comparando -se os resultados de acordo com as Grandes Regiões do país, a Categoria Administrativa e a Organização Acadêmica.
Considerando -se as notas médias dos alunos em cada Grande Região, observa -se que, em relação às notas aferidas em âmbito nacional ( média 37,7 para Ingressantes e 43,1 para Concluintes ), existe diferença estatisticamente significativa ao nível de 95 \%, para Ingressantes, entre a menor média, a obtida na região Norte ( 30,7 ) e a maior, na região Sudeste ( 39,4 ), nesse grupo.
Entre os Concluintes, também existe diferença estatisticamente significativa das médias ao nível de 95 \%.
A maior média de notas foi obtida pela região Norte ( 50,8 ), ao passo que a menor foi identificada na região Centro-Oeste ( 33,9 ).
As demais médias de notas foram 39,0 para a região Nordeste, 43,7 para a região Sudeste e 47,7 para a Sul.
Levando -se em conta os grupos dos estudantes em Categorias Administrativas, observa -se que a pontuação média mais elevada entre os Concluintes foi encontrada em instituições públicas ( 48,0 ), com desempenho acima da média nacional.
Já os Concluintes provenientes de instituições privadas obtiveram média 40,8, abaixo da média nacional.
Quanto aos Ingressantes, as notas mais elevadas, também acima da média nacional para esse grupo de estudantes, ocorreram nas instituições públicas ( 40,8 ).
Os Ingressantes vinculados a instituições privadas obtiveram nota média 34,0, abaixo da média nacional.
Tendo como foco as notas médias dos estudantes agrupados por Organização Acadêmica, distribuídas entre Universidades, Centros Universitários e Faculdades, para os Concluintes vinculados a Universidades, a média 45,2 ficou acima da média nacional.
Pode -se verificar a mesma realidade entre os Ingressantes de Universidades, cuja nota média foi 39,5, também acima da média nacional para esse grupo.
Considerando -se as demais Organizações, Concluintes de Centros Universitários e de Faculdades tiveram notas médias 41,9 e 39,2, respectivamente.
Ingressantes das mesmas Organizações Acadêmicas obtiveram médias 34,7 e 33,7, respectivamente. 3.1.2 Estatísticas Básicas no Componente de Formação Geral A Tabela 3.2 apresenta as estatísticas básicas em relação ao componente da prova que avalia a Formação Geral dos estudantes.
Os Concluintes obtiveram desempenho médio 43,9, que foi inferior ao dos Ingressantes, 44,4.
Quanto à variabilidade das notas nos dois grupos, o desvio padrão das notas obtidas foi 18,9 para Ingressantes e 21,0 para Concluintes.
Os dois grupos alcançaram notas máximas próximas, sendo 92,0 a maior nota obtida por pelo menos um Ingressante, e 95,0 a maior nota entre os Concluintes.
O Gráfico 3.3 propicia a comparação entre Concluintes e Ingressantes relativamente ao desempenho no componente de Formação Geral.
Dos Ingressantes, a faixa onde houve maior frequência deste grupo de estudantes foi entre 50 e 60 pontos, com 19,6 \% deles, sendo essa a faixa modal.
Esse também foi o intervalo em que houve concentração mais expressiva de Concluintes, 17,8 \%.
Nas faixas de notas seguintes, mais elevadas, houve uma alternância entre os dois grupos quanto a maior frequência.
Para os Concluintes, houve percentual maior de estudantes no intervalo entre 70 e 80 pontos, com 8,0 \% dos participantes, bem como entre 80 e 90, com 2,9 \%.
Na menor faixa de notas, que vai de zero a 10 pontos, Concluintes tiveram incidência superior a de Ingressantes : 9,1 \% para os primeiros e 5,0 \% para os segundos.
O coeficiente de assimetria foi negativo ( -0,2 ) para os dois grupos.
No Gráfico 3.4 são apresentadas as informações referentes ao desempenho dos Ingressantes e dos Concluintes no componente de Formação Geral, em diferentes agregações como Grande Região do país, Categoria Administrativa e Organização Acadêmica.
Considerando -se as notas médias dos estudantes segundo cada Grande Região, observa -se que no grupo de Ingressantes existe diferença estatisticamente significativa ao nível de 95 \% entre as médias na região Sudeste ( 46,1, a mais alta ) e Norte ( 35,9, a mais baixa ).
No grupo de Concluintes, a nota mais elevada ocorreu na região Sul, com média 48,0.
A diferença entre notas mais altas e mais baixas foi também estatisticamente significativa ao nível de 95 \% : a menor nota foi 26,0, correspondente aos Concluintes da região Centro-Oeste.
Destaca -se que Ingressantes superaram os Concluintes nas regiões Nordeste, Sudeste e Centro-Oeste, sendo que nessas duas últimas regiões, obtiveram médias superiores à nacional para Ingressantes.
Tendo como foco as Categorias Administrativas, as instituições públicas obtiveram médias mais elevadas do que as privadas : para Concluintes, notas médias de 50,3 em cursos de IES públicas e de 40,7 nos das privadas.
Para os Ingressantes, as notas foram de, respectivamente, 48,4 e 39,7.
Em comparação à média nacional, as notas de Ingressantes e de Concluintes vinculados às instituições públicas foram superiores, ao passo que, nas instituições privadas, foram inferiores.
Considerando -se o tipo de Organização Acadêmica, nota -se diferença estatisticamente significativa a 95 \% entre a maior e a menor média, tanto para Concluintes como para Ingressantes.
Nas Universidades, houve desempenho melhor do que em Centros Universitários e Faculdades.
Neste primeiro tipo de Organização Acadêmica, Ingressantes tiveram nota média 46,7, enquanto Concluintes, 46,3.
Nos Centros Universitários e Faculdades, os valores foram mais baixos : os Ingressantes tiveram média 39,0 nos Centros Universitários e 40,7 nas Faculdades ; entre Concluintes, a média foi 42,7 para Centros Universitários e 38,9 para as Faculdades. 3.1.3 Estatísticas Básicas do Componente de Conhecimento Específico A Tabela 3.3 apresenta as estatísticas básicas referentes à parte do componente de Conhecimento Específico da Área de Terapia Ocupacional.
A média do desempenho foi 38,7.
Os Concluintes apresentaram desempenho médio ( 42,9 ) superior ao dos Ingressantes ( 35,4 ) ; as medidas de dispersão – desvio padrão e amplitude geral das notas – são menores no grupo dos Ingressantes do que no dos Concluintes.
Os coeficientes de variação, 35,9 \% no grupo dos Ingressantes e 34,5 \% no dos Concluintes, indicam dispersão relativa menor nas notas desses últimos.
As notas máximas obtidas por Concluintes e Ingressantes foram iguais a 79,8 e 69,1, respectivamente, notas inferiores àquelas encontradas no componente de Formação Geral ( 95,0 e 92,0, respectivamente ).
As medianas das distribuições das notas dos Ingressantes e dos Concluintes foram 35,4 e 44,5, respectivamente, enfatizando a superioridade dos Concluintes neste componente da prova.
Assim como os Gráficos 3.1 e 3.3, o Gráfico 3.5, apresentado a seguir, proporciona uma comparação do desempenho entre Concluintes e Ingressantes, em relação ao componente de Conhecimento Específico.
Dentre as três comparações feitas, observa -se que, embora a distribuição das notas seja semelhante à apresentada na prova como um todo, neste componente, a diferença de desempenho entre Concluintes e Ingressantes é mais acentuada, destacando -se o desempenho mais elevado dos primeiros.
Claramente a distribuição das notas dos Ingressantes, mais concentradas à esquerda em relação aos Concluintes, evidencia o percentual mais alto de notas nas faixas mais baixas para este grupo.
A maior concentração de Concluintes está no intervalo ( 30 a 40 ] pontos com 32,0 \% das notas do grupo, ao passo que a de Ingressantes está no intervalo subseqüente ( 40 a 50 ] pontos com 29,4 \%.
Assim como os Gráficos 3.2 e 3.4, o Gráfico 3.6 apresenta uma comparação dos resultados em relação à Grande Região do país, à Categoria Administrativa e à Organização Acadêmica, agora levando em conta o desempenho dos Ingressantes e dos Concluintes no componente de Conhecimento Específico da prova.
Ao se observar as notas médias dos estudantes, segundo cada Grande Região do país, verifica -se que, entre os Concluintes, a mais elevada foi encontrada na região Norte ( 52,7 ), maior que a média nacional ( 42,9 ).
Em relação aos Ingressantes, a região Sul apresentou a média mais elevada ( 37,3 ), seguida da região Sudeste, com média igual a 37,1.
A média mais baixa entre os Concluintes ( 36,5 ) foi verificada na região Centro-Oeste, e entre os Ingressantes ( 28,9 ), na região Norte.
Quanto à Categoria Administrativa, observa -se comportamento semelhante àquele da parte relativa à Formação Geral e à prova como um todo, ou seja, as médias mais elevadas tanto entre os Concluintes quanto entre os Ingressantes foram encontradas em instituições de ensino públicas.
Nesse tipo de IES, a nota média de Concluintes foi 47,2, enquanto a de Ingressantes, 38,2.
Tais médias estão situadas acima da média nacional.
Nas instituições privadas, as médias foram, mais uma vez, inferiores em relação às públicas.
Concluintes obtiveram média 40,8 e Ingressantes, 32,1.
Em relação à Organização Acadêmica, o desempenho dos Concluintes foi superior ao dos Ingressantes nas três categorias.
Nas Universidades, os Concluintes obtiveram a média mais elevada, igual a 44,8.
Um pouco abaixo estiveram os Centros Universitários ( 41,6 ) e as 3.2 ANÁLISE DAS QUESTÕES OBJETIVAS 3.2.1 Componente de Formação Geral A Tabela 3.4 apresenta as estatísticas básicas relativas às oito questões objetivas do componente da prova que abrange a Formação Geral dos estudantes.
Como pode ser observado, os estudantes Concluintes obtiveram desempenho superior, com nota média de 50,7, em relação aos alunos Ingressantes, cuja nota média foi 49,6.
As medidas de dispersão nos dois grupos de alunos foram similares, como indicam os desvios padrões, assim como a amplitude total das notas, sendo que a nota mínima foi zero e a máxima, 100,0, para os dois grupos de alunos avaliados.
Os coeficientes de variação, 42,7 \% no grupo dos Ingressantes e 46,4 \% no dos Concluintes, indicam maior dispersão relativa nas notas desses últimos.
A mediana foi igual a 57,1 em ambos os grupos.
O coeficiente de assimetria do total de estudantes revelou assimetria negativa ( -0,3 ).
A Tabela 3.5 apresenta as classificações das questões objetivas do componente de Formação Geral, segundo o índice de facilidade.
As questões objetivas da prova foram assim avaliadas : das oito questões, nenhuma teve o índice de facilidade classificado como muito fácil.
Três questões foram tidas como fáceis, por terem índice de acertos situado na faixa entre 0,61 e 0,85 ( de 61,0 \% a 85,0 \% de acertos ).
Duas questões foram consideradas de dificuldade média, situando -se no intervalo entre 0,41 e 0,60 do índice de facilidade, ou seja, houve entre 41,0 \% e 60,0 \% de acertos.
Duas questões foram classificadas como difíceis, situando -se no intervalo entre 0,16 e 0,40.
Por fim, uma questão apresentou menos de 0,15 de acertos, razão pela qual foi classificada como muito difícil.
A Tabela 3.6 descreve os resultados das análises das questões objetivas relativas à Formação Geral, segundo o poder de discriminação, utilizando -se para tal, do índice de discriminação ponto bisserial.
Cinco das oito questões apresentaram índices acima de 0,40 e, assim, foram classificadas como muito boas nesta característica para esse grupo de alunos ; duas questões tiveram bom índice de discriminação, entre 0,30 e 0,39.
Nenhuma questão foi classificada como média.
A questão de número 6 teve fraco nível de discriminação e foi eliminada do cômputo da nota final.
Na sequência, na Tabela 3.7, estão especificadas, por questão, os resultados da análise relativa aos índices de facilidade e de discriminação de cada questão no componente relativo à Formação Geral da prova de Terapia Ocupacional.
O índice de facilidade variou de 0,08 a 0,75, e o de discriminação, de 0,16 a 0,56.
As questões com índices de discriminação muito bom, de números 1, 3, 4, 5 e 8, figuraram entre as mais fáceis desse conjunto : três classificada pelo índice de facilidade na categoria fácil ( questões 3, 4 e 8 ) e duas na categoria médio ( questões 1 e 5 ).
A questão 8 foi a que apresentou maior poder discriminatório, com índice 0,56, e foi também a mais fácil, com uma proporção de 0,75 de acertos.
A questão de número 6 apresentou índice de discriminação 0,08, ou seja, um quantitativo pequeno de estudantes conseguiu resolvê -la, dentro do universo de participantes.
Além disso obteve índice de discriminação fraco, o que a eliminou do cômputo final das notas.
A questão 7 obteve índice de discriminação bom, 0,35.
O Gráfico 3.7 analisa o comportamento da questão de número 8 de Formação Geral.
Trata -se da questão mais fácil e com o maior índice de discriminação desse componente da prova.
No gráfico a seguir, cada uma das cinco curvas representa o percentual de respostas em determinada alternativa da questão, em função da nota dos estudantes nesta parte da prova ( Formação Geral/Múltipla Escolha ) antes de possíveis eliminações pelo critério do ponto bisserial.
A curva em rosa corresponde à alternativa B, a correta para este item.
Assim, observa -se que entre os estudantes com notas mais baixas, com desempenho mais fraco nesta parte da prova, a situação mais frequente foi a escolha de uma das alternativas incorretas : a alternativa C ( em verde ) ou E ( em roxo ).
Na medida em que a nota aumenta, indicando desempenho melhor nesta parte da prova, aumenta concomitantemente a proporção de estudantes que selecionou a alternativa correta B, atingindo 100 \% para as notas acima de 85,0.
Essa análise permite verificar como a questão discriminou entre os grupos de desempenho, justificando o alto índice obtido na questão.
Os gráficos relativos às demais questões de Formação Geral constam do Anexo I. 3.2.2 Componente de Conhecimento Específico A Tabela 3.8 apresenta as estatísticas básicas em relação às questões objetivas do componente de Conhecimento Específico da prova.
Nesse conjunto de questões, os estudantes Concluintes ( 47,1 ) obtiveram nota média superior aos Ingressantes ( 40,0 ).
Quanto à variabilidade das notas, o desvio padrão no grupo dos Ingressantes foi 14,2 e no grupo dos Concluintes, 16,0.
Os coeficientes de variação foram 35,5 \% e 34,0 \%, respectivamente, indicando menor dispersão relativa nas notas dos Concluintes.
Com relação à mediana, no grupo dos Ingressantes, esta foi igual a 41,7, enquanto no grupo dos Concluintes, a mediana foi igual a 50,0, ambas superando um pouco a média do grupo correspondente.
O coeficiente de assimetria foi negativo Ingressantes ( -0,4 ) e para Concluintes ( -0,8 ).
A Tabela 3.9 apresenta as classificações das questões objetivas da parte da prova relativa ao componente de Conhecimento Específico, segundo o índice de facilidade.
Deve -se considerar o fato de a questão 21 ter sido anulada pela Comissão, de modo que a classificação foi estabelecida com base em 26 das 27 questões.
A partir dos índices obtidos, pode -se concluir que um pouco mais da metade das questões objetivas da prova foram tidas como pelo menos difíceis : das 26 questões, doze foram classificadas como difíceis e duas como muito difíceis.
Não houve questão classificada como muito fácil, ao passo que cinco foram tidas como fáceis, na faixa de 0,61 a 0,85 do índice de facilidade, e sete consideradas médias, entre 0,41 e 0,60.
Na Tabela 3.10 são apresentados os índices de discriminação das questões objetivas do componente de Conhecimento Específico da prova.
Observa -se, segundo os índices estabelecidos, que a capacidade de discriminação das questões objetivas desse componente da prova foi classificada, para sete das 26 questões válidas, como boa, enquanto seis delas tiveram índice de discriminação muito bom.
A metade das questões – 13 em 26 – proporcionou, portanto, índices de discriminação bons ou muito bons.
Onze A Tabela 3.9 apresenta as classificações das questões objetivas da parte da prova relativa ao componente de Conhecimento Específico, segundo o índice de facilidade.
Deve -se considerar o fato de a questão 21 ter sido anulada pela Comissão, de modo que a classificação foi estabelecida com base em 26 das 27 questões.
A partir dos índices obtidos, pode -se concluir que um pouco mais da metade das questões objetivas da prova foram tidas como pelo menos difíceis : das 26 questões, doze foram classificadas como difíceis e duas como muito difíceis.
Não houve questão classificada como muito fácil, ao passo que cinco foram tidas como fáceis, na faixa de 0,61 a 0,85 do índice de facilidade, e sete consideradas médias, entre 0,41 e 0,60.
Na Tabela 3.10 são apresentados os índices de discriminação das questões objetivas do componente de Conhecimento Específico da prova.
Observa -se, segundo os índices estabelecidos, que a capacidade de discriminação das questões objetivas desse componente da prova foi classificada, para sete das 26 questões válidas, como boa, enquanto seis delas tiveram índice de discriminação muito bom.
A metade das questões – 13 em 26 – proporcionou, portanto, índices de discriminação bons ou muito bons.
Onze delas foram classificadas como médias e outras duas como fracas, sendo treze, por conseguinte, a quantidade de questões nos dois patamares mais baixos de discriminação.
Constata -se, assim, que a prova – no que se refere ao componente de Conhecimento Específico – possuía capacidade mediana de discriminar entre aqueles que dominam ou não o conteúdo.
As questões 11 e 27, com fraco índice de discriminação foram eliminadas do cômputo das notas finais.
Na Tabela 3.11 estão especificados, por questão, os resultados da análise relativa aos índices de facilidade e de discriminação de cada uma das questões do componente de Conhecimento Específico da prova de Terapia Ocupacional.
Dentre as questões que alcançaram os maiores índices de discriminação, muito bom, as de números 15, 22, 24 e 25 foram classificadas como fáceis pelo índice de facilidade e as questões 28 e 35 como médias.
Em termos de proporção, tais questões variaram de 0,47 a 0,73 quanto aos acertos.
Seus índices de discriminação situaram -se de 0,41 a 0,47.
A questão de número 27 foi a mais difícil dentre as 26 questões específicas válidas, com baixo índice de facilidade, 0,08, o que, em termos percentuais, corresponde a apenas 8,0 \% de acertos.
Essa questão apresentou poder discriminatório igualmente baixo, 0,11, o que comprova ter sido esta a mais difícil para os estudantes e como já comentado foi eliminada pelo critério do índice de discriminação.
Destaca -se, também, a questão de número 11, para a qual o índice de facilidade foi 0,20, o que, em termos percentuais, corresponde a 20,0 \% de estudantes que responderam acertadamente.
Foi também uma das mais difíceis da prova.
Obteve índice de discriminação fraco, 0,12 e foi eliminada por este critério.
O Gráfico 3.8 analisa o comportamento da questão 15 do componente de Conhecimento Específico.
Esta foi uma das questões mais fáceis da prova, apresentando índice de facilidade 0,66, ou seja, 66,0 \% dos estudantes assinalaram acertadamente a opção C, correspondente ao gabarito.
Seu índice de discriminação foi igual a 0,45, classificado como muito bom, embora outras questões que apresentaram percentuais mais elevados de respostas corretas tenham tido índices de discriminação igualmente altos.
Neste gráfico, cada uma das cinco curvas representa o percentual de respostas em determinada alternativa da questão 15, em função da nota dos estudantes nesta parte da prova antes de possíveis eliminações de questões pelo critério do ponto bisserial.
A alternativa correta C, representada no gráfico pela curva em verde, foi escolhida em maiores proporções pelos alunos com desempenho melhor nesse componente do exame, enquanto as alternativas incorretas, também denominadas distratores, foram selecionadas principalmente por aqueles com notas mais baixas, ainda que aqueles com nota zero, na sua quase totalidade tenham deixado em branco ou marcado mais de uma alternativa, comportamento considerado inválido.
A proporção de alunos que selecionaram a resposta correta C \textbf{cresce}, \textbf{grosso} modo, com a nota dos alunos chegando a atingir 100 \% para notas acima de 60,0, enquanto a proporção dos que escolhem alternativas incorretas, decai como função da nota nesta parte da prova.
Os gráficos relativos às demais questões do componente de Conhecimento Específico constam do Anexo I. 3.3 ANÁLISE DAS QUESTÕES DISCURSIVAS 3.3.1 Componente de Formação Geral A análise dos resultados de desempenho dos estudantes de Terapia Ocupacional, nas duas questões discursivas relativas à Formação Geral, encontra -se na Tabela 3.12 e no Gráfico 3.9.
Observa -se, a partir destes, que as notas médias foram mais baixas nesse conjunto do que no conjunto das objetivas.
A média geral dos Ingressantes nas questões objetivas de Formação Geral foi 49,6 ( Tabela 3.4 ), ao passo que nas questões discursivas, a média foi 36,5 para esse grupo.
O mesmo ocorreu entre os Concluintes, que tiveram, em Formação Geral, média 50,7 nas questões objetivas e 33,6 nas questões discursivas.
Cabe destacar que, nas questões discursivas, o desempenho dos Concluintes foi um pouco inferior ao dos Ingressantes.
Os dois grupos atingiram a mesma nota máxima, 100,0, o que acompanha a constatação de que não foi muito evidente uma possível superioridade dos Concluintes, em relação aos Ingressantes, nesta parte do exame.
Os coeficientes de assimetria indicam assimetria positiva para os Concluintes, e comportamento simétrico para os Ingressantes.
O grupo modal de notas foi o intervalo [ 0 ; 10 ] para ambos os grupos.
Verifica -se a semelhança das distribuições das notas desses grupos, no Gráfico 3.9, apresentado a seguir.
Na sequência, os resultados verificados para cada uma das questões discursivas de Formação Geral serão apresentados, estabelecendo -se relações com os conteúdos englobados em cada uma delas.
Os comentários da banca corretora a respeito do observado nas respostas dos estudantes, suas impressões e conclusões serão apresentados junto à análise de cada questão.
Cumpre esclarecer que, tendo em vista que as questões discursivas de Formação Geral são padronizadas, ou seja, constam de todas as provas, os comentários da banca são os mesmos para todas as carreiras acadêmicas, sendo direcionados a todos os estudantes que participaram do ENADE/2010.
A seguir, serão analisados os desempenhos de Ingressantes e Concluintes da Área de Terapia Ocupacional nas duas questões discursivas de Formação Geral do ENADE/2010. 3.3.1.1 Análise da Questão Discursiva 9 do Componente de Formação Geral Os dados obtidos a partir das respostas à questão 9 encontram -se na Tabela 3.13 e no Gráfico 3.10.
Nessa questão – de melhor desempenho dentre as duas de Formação Geral – os Ingressantes tiveram média 39,4, e os Concluintes, 36,6, resultando com diferença de 2,8 pontos entre os dois grupos, favorável aos Ingressantes.
Houve semelhança na distribuição de notas dos dois grupos, observando -se percentuais próximos de Concluintes ( 32,2 \% ) e Ingressantes ( 28,1 \% ) que deixaram a questão em branco, bem como daqueles que mesmo não tendo deixado a questão em branco obtiveram nota zero na mesma, com 3,7 \% e 2,6 \%, respectivamente.
Entre Ingressantes, a nota mediana foi 45,0 pontos, enquanto Concluintes obtiveram 40,0.
Considerando -se somente as notas não nulas, Ingressantes apresentaram como faixa modal os intervalos ( 40 a 50 ] e ( 50 a 60 ] com 12,3 \% das notas em cada, ao passo que para os Concluintes as faixas modais foram ( 50 a 60 ] e ( 60 a 70 ] com 11,5 \% cada.
A nota máxima, 100,0, foi alcançada tanto por Concluintes quanto por Ingressantes. 3.3.1.2 Comentários sobre as respostas à Questão Discursiva 9 No que tange mais precisamente à atualidade do tema proposto e ao desenvolvimento do conteúdo, deve ser ressaltado que se mostrou sintonizado com a realidade política brasileira, bem como com a necessária reflexão sobre o papel na sociedade do profissional formado pelas Instituições de Ensino Superior do país.
Dissertar sobre democracia e ética foi algo apropriado para a época em que a prova foi aplicada, pois coincidiu com o período das eleições, em que tais temas são mais discutidos na vida cotidiana.
Não se pode, entretanto, deixar de se levar em conta a tendência de os estudantes terem limitado a associação do tema “ ética ” ao assunto “ política ”, considerando -se que a ética abarca não apenas a política.
Tal associação atingiu, em alguma medida, o alcance mais preciso das respostas.
As respostas apresentadas refletiram motivação para escrever sobre o tema da questão.
Alguns estudantes chegaram a desenvolver textos eloquentes e engajados.
Em alguns casos, o conceito de democracia, por exemplo, foi bem desenvolvido ao longo do texto, simultaneamente às exemplificações pedidas.
Além disso, muitos estudantes associaram o exercício da ética ao seu futuro ofício, reconhecendo o seu papel como cidadão no desempenho de sua profissão.
Essa associação insinuou, também, o esforço de reconhecimento da noção de ética prática, ainda que de maneira codificada.
Um dos pontos positivos revelou -se no fato de, na maioria das respostas analisadas, haver uma preocupação em definir a sociedade democrática.
Isso facilitou a correção e elevou a média de notas.
Outro aspecto positivo foi a presença pertinente de exposições coerentes de comportamento ético de um futuro profissional.
A articulação dos itens “ a ” e “ b ” da questão produziu, portanto, todo tipo de reflexão acerca de uma postura ética na sociedade, inclusive, corporativa.
A comprovação de que o tema ética tem sido abordado nos cursos superiores foi outro destaque positivo detectado durante as correções.
Assim, além do caráter crítico de que muitas das respostas corretas se revestiram, os relatos constantes sobre desdobramentos da ética comprovaram a abordagem sobre o tema, no mínimo, a ética profissional, nas carreiras universitárias que participaram do ENADE.
A despeito de alguns estudantes se posicionarem contra o exame, seja anulando a questão ou escrevendo algum comentário desabonador sobre o processo, foi observado que significativa parte dos participantes tem consciência político-acadêmica e percebe por que está realizando aquele tipo de prova.
Desse modo, a maioria dos estudantes tinha algo a “ dizer “ sobre o tema da questão, isto é, valia -se da oportunidade para se posicionar a respeito do assunto trazido à baila.
Ainda do ponto de vista do conteúdo desenvolvido nas respostas, observou -se preocupação solidária e consciência cidadã dos estudantes.
Mais um destaque positivo crucial se refere à habilidade textual.
Os estudantes, por exemplo, expressaram atitudes éticas e não éticas e as contrapunham ao conceito de democracia.
Há casos em que os alunos definiam a corrupção política como um exemplo de atitude não ética, enfatizando a real missão desse cargo em uma sociedade democrática : “ atender aos anseios da população e zelar pelo bem comum ”, escreveram vários estudantes.
No tocante à produção de texto dissertativo, à objetividade e ao caráter sintético, muitos estudantes corresponderam ao solicitado, elaborando pequenos textos dissertativos.
Nessas respostas, percebeu -se demonstração de domínio do assunto focado e utilização dos textos de referência como base para a elaboração.
O enunciado da questão 9 determinou que fosse elaborado um texto dissertativo que contivesse a )..., b )... e c )..., fornecendo um encaminhamento para as respostas que conduziu tanto a um resultado positivo quanto a um negativo, a saber : foi produtiva a resposta dos estudantes no sentido de que conseguiram, na maior parte das vezes, atender plenamente aos três quesitos apresentados na questão ; por outro lado, mais especificamente devido ao uso das três letras indicativas ( a, b e c ), houve tendência à elaboração de respostas divididas em itens, em vez de um texto eminentemente dissertativo.
Muitos estudantes apresentaram respostas nesse formato fragmentado para satisfazer a formalidade do roteiro sugerido no comando da questão.
A ideia-chave que deveria permear toda a resposta – democracia – se perdia em fragmentos desconexos entre si.
Se, por um lado, o roteiro do comando da questão facilitou o raciocínio de alguns dos estudantes, por outro, induziu grande parte deles a dois caminhos : i ) o das respostas curtas e objetivas, do tipo “ um profissional de saúde não ético é aquele que conta para as outras pessoas o que foi feito na consulta com o seu paciente ” e ii ) o outro caminho, que é do contraponto entre o conceito de sociedade democrática e a realidade brasileira.
Esse segundo caminho levou vários universitários a produzirem textos bem escritos, distantes, entretanto, do tema exigido.
No que se reporta às especificidades do desenvolvimento do conteúdo, a dificuldade ou falta de competência comunicativa revelada em grande parte das respostas denuncia a pouca exposição dos alunos à leitura de textos escritos de mais \textbf{excelência}.
Tal dificuldade mostrou -se mais evidente na impropriedade no uso de conectores e na incapacidade de elaboração de texto dissertativo, no tratamento de questões mais abstratas.
Muitos estudantes demonstraram não saber em que consistia a ética e confundiram exemplificação com definição : ao tentarem dar exemplos, via de regra, elaboravam uma definição incipiente do que se pedia.
Portanto, a dificuldade dos estudantes em argumentar e articular as ideias com coesão e coerência ainda são os grandes “ vilões ” desse processo \textbf{avaliativo}.
O excesso de fragmentos e textos prévios acerca do tema permitiu que a elaboração das respostas se resumisse a uma colagem ou paráfrase mal feita dos materiais disponíveis.
Em vez de subsidiar as respostas, esse material acabou por influenciar negativamente o desempenho dos estudantes.
Desse modo, verificou -se um nível de letramento ainda significativamente baixo, visto que a maioria apresentou uma visão bastante reducionista e simplificada sobre o tema em pauta.
De um modo geral, ao conceituarem ética e sociedade democrática, os estudantes demonstraram desconhecer o limite entre os dois conceitos, sem aludirem à \textbf{interseção} de ambos.
Foi expressiva a quantidade de redações nas quais questionavam a democracia pelo simples fato de estarem fazendo, a fortiori, aquela prova especificamente.
Em muitas das respostas, os estudantes se diziam incomodados pelo fato de estarem obrigados a realizar o exame, demonstrando sua insatisfação.
Por diversas vezes, os estudantes não foram capazes de sintetizar e de selecionar seu discurso o suficiente para atender ao que era pedido, o que ocasionou a falta de, pelo menos, um dos aspectos da resposta.
Foi constatada grande parte de respostas com problemas graves em termos textuais e gramaticais, em relação ao que se esperaria de pessoas que estão vinculadas ao Ensino ( Técnico ) Superior.
Alguns estudantes sequer conseguiam elaborar uma frase coerente e coesa, apresentando respostas que poderiam ser classificadas como de analfabetismo funcional.
De maneira ampla, manifestou -se da parte dos estudantes uma grande dificuldade no trato da língua escrita, dificuldade essa que se traduz, sobretudo, por um desrespeito bastante frequente às regras relativas : i ) à ortografia ( por exemplo : “ proficionais ”, “ infligir as leis ”, “ corrupitos ”, “ democrácia ”, “ opitar ”, “ manter o cizilo proficionau ” ), troca de “ sociedade democrática ” por “ sociedade demográfica ” e de “ ética ” por “ étnica ” ; ii ) à pontuação ; iii ) à concordância verbal ( “ os político deve (... ) eles pode ” ) ; iv ) ao modo de organização textual ( sem distinção entre o que seja dissertativo, argumentativo ou narrativo ).
Frisa -se que esse comentário não se restringe ao que se costuma chamar de “ língua padrão ”, mas ao que poderia ser chamado de “ \textbf{excelência} textual ”, que não se mede exclusiva ou enfaticamente pela correção gramatical, mas pela competência, objetividade, clareza e, em alguma medida, pela criatividade no ato de redigir.
Por fim, destaca -se um dos problemas dignos de nota na formulação da questão.
A transcrição do Dicionário Houaiss apresentava defeitos.
Por exemplo, a definição 5 de “ democracia ” indicava que esse é o tipo de “ governo que acata a vontade da maioria, embora respeitando os direitos e a livre expressão das minorias ”.
O emprego de um valor concessivo na oração adverbial, por oposição a um valor aditivo, contradiz as concepções efetivamente democráticas apresentadas de 1 a 4.
Repercussões desse valor concessivo, mal empregado, foram encontradas em incontáveis respostas, as quais diziam ser típico de um governo democrático, seguir a vontade da maioria, “ apesar ” de respeitar as minorias.
Outro desses problemas dignos de nota refere -se à indicação do dicionário quanto a “ governo ”, ao passo que o comando focava a ideia mais ampla de “ sociedade democrática ”. diziam ser típico de um governo democrático, seguir a vontade da maioria, “ apesar ” de respeitar as minorias.
Outro desses problemas dignos de nota refere -se à indicação do dicionário quanto a “ governo ”, ao passo que o comando focava a ideia mais ampla de “ sociedade democrática ”. 3.3.1.3 Análise da Questão Discursiva 10 do Componente de Formação Geral O desempenho dos estudantes na questão 10 foi inferior ao obtido na de número 9, com nota média 33,5 para Ingressantes e 30,5 para Concluintes.
A nota máxima também foi atingida em ambos os grupos, sendo que os percentuais de alunos que deixaram a questão em branco, tanto no grupo dos Ingressantes ( 32,3 \% ) quanto no dos Concluintes ( 37,8 \% ), foram superiores àqueles apresentados na questão 9, bem como a proporção de alunos com nota zero, mesmo não tendo deixado a questão em branco.
Houve maior percentual de alunos, considerando -se a distribuição de notas não nulas, no intervalo entre 40 e 50 pontos, tanto para Ingressantes quanto para Concluintes, respectivamente 14,4 e 14,8 \%.
Não se manifestou uma superioridade de Concluintes nas diferentes faixas de notas posteriores a esse intervalo, com pequena superioridade de Ingressantes no desempenho.
Essas informações estão ilustradas na Tabela 3.14 e no Gráfico 3.11. 3.3.1.4 Comentários sobre as respostas à Questão Discursiva 10 Inicialmente, será apresentada uma panorâmica sobre as respostas, com destaque para os aspectos positivos observados.
Em relação ao desenvolvimento do tema e à vinculação do mesmo à vida cotidiana dos estudantes que participaram do exame, a visão positiva em relação à Educação de Jovens e Adultos ( EJA ) reforça a ideia de que a educação pode trazer crescimento não somente profissional, mas, sobretudo, para ampliar os horizontes da cidadania.
Foram frequentes as respostas que apontavam a educação como a melhor saída, senão a única, para o resgate da cidadania.
A atuação dos governos, em especial o governo federal, foi um ponto destacado nas respostas, as quais sinalizavam para a necessidade de as instâncias governamentais assumirem um papel cada vez mais ativo na consolidação da educação – pública, sobretudo – no país.
As políticas públicas focadas na educação foram, então, um item bastante frequente nas respostas, seja para enaltecer, criticar ou requisitar a ação governamental.
Outro ponto positivo destacado diz respeito à consciência ou ao reconhecimento da precariedade de condições de trabalho dos professores, em geral, no Brasil e, em particular, daqueles ocupados com turmas de jovens e adultos.
Nesse sentido, foram inúmeras as propostas apresentadas na segunda parte das respostas com relação à imperiosidade de melhoria salarial para os docentes, e à melhoria das condições gerais de trabalho desses profissionais.
A associação direta entre valorização do professor e sucesso escolar dos estudantes foi um ponto complementar a essa preocupação com os docentes.
Os estudantes demonstraram uma visão de vínculo no processo de ensino-aprendizagem, no qual alunos e professores são partes inalienáveis entre si.
A questão da repetência e da evasão escolar, referenciadas no suporte da questão, foram então, abordadas pelos alunos por meio do viés dessa valorização docente.
Ainda que os temas específicos da repetência e da evasão não fossem contemplados a contento na maior parte das respostas, ao menos, nessa associação com a qualificação e a valorização dos professores, algo foi trabalhado analiticamente pelos estudantes.
Num plano mais abrangente, houve reconhecimento de que, no processo ensino-aprendizagem, não há atores isolados : alunos, profissionais de educação, governos e a própria sociedade como corresponsáveis pelo sucesso ou fracasso de tal processo.
A atualidade do tema e a sua explícita vinculação com o cotidiano de muitos estudantes que realizavam o exame fizeram com que os respondentes demonstrassem um visível grau de envolvimento com o tema proposto.
Ou seja, foi possível detectar um empenho notório dos estudantes ao redigir suas respostas com seriedade e comprometimento relativo à delicada problemática dos milhões de estudantes trabalhadores do país.
Utilizar todo o espaço destinado à resposta, procurando aproveitar ao máximo a oportunidade de emitir um argumento ou uma opinião revelou -se como uma das evidências do aludido empenho.
A EJA também foi reconhecida como um dos caminhos concretos para o acesso à Universidade, no médio prazo.
As respostas se voltaram, muitas vezes, à condição de depoimentos pessoais sobre a conciliação de trabalho e estudo.
Nessa direção, ao assumirem o papel de depoentes, os estudantes expuseram suas histórias de vida, suas expectativas e percalços, suas esperanças e entraves cotidianos, seus anseios mais ambiciosos e problemas estruturais, que poderiam se resolver por meio da educação.
Nesse inciso, ficava clara a relação que se formava entre educação e mobilidade social, mediatizada por iniciativas governamentais como a EJA.
Uma sugestão plausível e válida para a solução dos problemas com a EJA no país reportou -se à necessidade de maior participação do público alvo na definição de estratégias para essa modalidade de ensino.
Nesse detalhe, era notória a associação das preocupações relativas à questão 9, que trata do problema da democracia e da cidadania, com o problema específico da educação escolar.
Houve um número significativo de sugestões ou propostas para se combater o problema da repetência e da evasão, destacando -se melhores salários para os professores e aulas mais dinâmicas.
A lista dessas sugestões, contudo, era bem mais ampla e diversificada, incluindo ainda : a redução da jornada de trabalho dos jovens e adultos que participam da EJA ; melhoria de infraestrutura nas instituições responsáveis pela EJA ; incentivo financeiro aos estudantes trabalhadores por meio de um sistema de bolsas de estudo ; realização de campanhas motivacionais para o retorno aos estudos ; estímulo à participação das empresas, por meio da solicitação de recursos materiais, espaço físico para as salas, auxílio financeiro ; incremento dos cursos profissionalizantes no país ; e ampliação da oferta de ensino a distância em todo o território nacional.
Muitos estudantes abordaram, com propriedade, a associação entre um mercado de trabalho cada vez mais competitivo e a necessidade de qualificação profissional, inclusive com uma clara percepção de que os salários são tanto maiores quanto maior a quantidade de anos de estudo e/ou a qualidade da formação.
Nesse tópico, retorna -se ao papel crucial da EJA como viabilizadora dessa formação.
Em continuidade à análise, serão apresentados aspectos negativos que foram notados com frequência nas respostas à questão 10.
De um ponto de vista mais global, a dificuldade repousa no entendimento da questão, isto é, na decifração do comando.
Muitos fugiram completamente do assunto proposto, ainda que o enunciado da questão fosse farto em dados específicos ao tema.
Em vez de responderem diretamente quais eram os problemas enfrentados por aqueles que têm de trabalhar e estudar ao mesmo tempo, muitos estudantes optavam por falar das dificuldades que existem na educação brasileira, de forma estendida e pasteurizada.
Vários estudantes não conseguiam, em suas respostas, relacionar as dificuldades de conciliação da escola com o trabalho, a não ser, como já mencionado, nos casos de depoimentos pessoais, relatos de experiência, ou seja, traços de uma biografia.
Assim, eram respostas cujo teor era apenas o de registrar o abandono dos estudos por parte de vários jovens que têm de trabalhar.
O tom das respostas, nesses casos, era meramente de uma descrição e não de uma análise, muito menos de um elemento que compusesse uma argumentação.
A dificuldade de compreensão da questão se expressava nas respostas que, em lugar de mostrarem os problemas de conciliação entre escola e trabalho, prendiam -se à ideia de “ associação ” – termo que consta do enunciado da questão – entre escola e trabalho.
Muitos alunos relatavam as consequências positivas de uma boa educação no momento em que se busca bons postos de trabalho.
Tratava -se de uma fuga parcial ao comando da questão, evidenciando má compreensão do mesmo.
Uma dificuldade facilmente percebida nas respostas refere -se à tendência de os estudantes tratarem dos problemas relativos à escola e, muito separadamente, daqueles problemas relativos ao trabalho, reconhecendo uma espécie de sinergia que impediria de visualizá -los em sua individualidade, ao mesmo tempo que em sua complementaridade. ampla e diversificada, incluindo ainda : a redução da jornada de trabalho dos jovens e adultos que participam da EJA ; melhoria de infraestrutura nas instituições responsáveis pela EJA ; incentivo financeiro aos estudantes trabalhadores por meio de um sistema de bolsas de estudo ; realização de campanhas motivacionais para o retorno aos estudos ; estímulo à participação das empresas, por meio da solicitação de recursos materiais, espaço físico para as salas, auxílio financeiro ; incremento dos cursos profissionalizantes no país ; e ampliação da oferta de ensino a distância em todo o território nacional.
Muitos estudantes abordaram, com propriedade, a associação entre um mercado de trabalho cada vez mais competitivo e a necessidade de qualificação profissional, inclusive com uma clara percepção de que os salários são tanto maiores quanto maior a quantidade de anos de estudo e/ou a qualidade da formação.
Nesse tópico, retorna -se ao papel crucial da EJA como viabilizadora dessa formação.
Em continuidade à análise, serão apresentados aspectos negativos que foram notados com frequência nas respostas à questão 10.
De um ponto de vista mais global, a dificuldade repousa no entendimento da questão, isto é, na decifração do comando.
Muitos fugiram completamente do assunto proposto, ainda que o enunciado da questão fosse farto em dados específicos ao tema.
Em vez de responderem diretamente quais eram os problemas enfrentados por aqueles que têm de trabalhar e estudar ao mesmo tempo, muitos estudantes optavam por falar das dificuldades que existem na educação brasileira, de forma estendida e pasteurizada.
Vários estudantes não conseguiam, em suas respostas, relacionar as dificuldades de conciliação da escola com o trabalho, a não ser, como já mencionado, nos casos de depoimentos pessoais, relatos de experiência, ou seja, traços de uma biografia.
Assim, eram respostas cujo teor era apenas o de registrar o abandono dos estudos por parte de vários jovens que têm de trabalhar.
O tom das respostas, nesses casos, era meramente de uma descrição e não de uma análise, muito menos de um elemento que compusesse uma argumentação.
A dificuldade de compreensão da questão se expressava nas respostas que, em lugar de mostrarem os problemas de conciliação entre escola e trabalho, prendiam -se à ideia de “ associação ” – termo que consta do enunciado da questão – entre escola e trabalho.
Muitos alunos relatavam as consequências positivas de uma boa educação no momento em que se busca bons postos de trabalho.
Tratava -se de uma fuga parcial ao comando da questão, evidenciando má compreensão do mesmo.
Uma dificuldade facilmente percebida nas respostas refere -se à tendência de os estudantes tratarem dos problemas relativos à escola e, muito separadamente, daqueles problemas relativos ao trabalho, reconhecendo uma espécie de sinergia que impediria de visualizá -los em sua individualidade, ao mesmo tempo que em sua complementaridade.
Isso significa que os estudantes distinguiam por mero isolamento, ou seja, eles não conseguiam distinguir e articular e, sobretudo, reconhecer as convocações mútuas entre os problemas da escola e do trabalho, no caso explícito dos estudantes trabalhadores.
Essa dificuldade de identificar, associar e analisar aspectos distintos, até mesmo concorrentes, antagônicos e complementares, denuncia a fragilidade de um pensamento simplificador que a qualidade das respostas revelou ser ainda hegemônico nas universidades brasileiras.
Notou -se a presença de preconceitos na análise dos fatores que levam ao abandono dos estudos.
Muitas respostas apontavam como causa desse abandono e da repetência a preguiça, a falta de perseverança e o cansaço dos estudantes.
Nesse tópico específico, fica evidente a responsabilização do indivíduo pelos percalços e fracassos de sua vida, descontextualizando -o de uma sociedade que, em seu movimento de conjunto, lhe impõe condicionantes.
A precariedade do uso da língua pátria de modo culto foi uma constante, conforme mencionado na análise da questão 9.
As impropriedades lexicais e a falta de elementos coesores que articulassem as partes constitutivas das respostas foram um incômodo constante.
Nesse sentido, agrava -se a situação da baixa qualidade textual em decorrência complementar da construção de textos meramente descritivos e não argumentativos, como é exigido no suporte da questão.
Isso demonstra a dificuldade dos estudantes universitários na compreensão de que uma argumentação requer, imperiosamente, justificações para os apontamentos apresentados.
Os textos das respostas eram, em sua maioria, opinativos, sem qualquer compromisso com elementos explicativos ou com a busca de justificativas. 3.3.2 Componente de Conhecimento Específico Na parte da prova relativa ao componente de Conhecimento Específico, observa -se que a diferença de desempenho foi mais acentuada entre Ingressantes e Concluintes e a favor dos Concluintes.
A Tabela 3.15 contém dados elucidativos, na qual é possível observar que, no grupo dos Ingressantes, a nota média 9,6 é inferior à alcançada por Concluintes, 19,0.
Essas médias foram menores do que as obtidas no conjunto das questões objetivas, nas quais Ingressantes alcançaram nota média 40,0, e Concluintes, 47,1 ( Tabela 3.8 ).
O Gráfico 3.12 apresenta a distribuição de notas nos diferentes intervalos.
Observa -se que tanto para Ingressantes quanto para Concluintes a faixa modal está entre zero e 10 pontos, com respectivamente 60,0 e 36,2 \% dos estudantes do grupo correspondente.
Para todas as faixas acima desta, a proporção de Concluintes é maior do que a de Ingressantes.
A análise de cada uma destas questões será feita a seguir. ∗ 3.3.2.1 Análise da Questão Discursiva 38 do Componente de Conhecimento Específico Na questão 38, cujos resultados aferidos encontram -se descritos na Tabela 3.16, o grupo dos Ingressantes obteve nota média 3,5, enquanto a média dos Concluintes foi 7,2. ∗ A análise completa, com apresentação de trechos de respostas de estudantes, pode ser encontrada no documento “ Relatório da Correção das Questões Discursivas ”, elaborado pelos docentes que participaram do processo de correção das mesmas.
Deixaram a questão em branco 66,6 \% dos Ingressantes e 59,3 \% dos Concluintes.
Além disso, as proporções de nota nula mesmo não tendo deixado a questão em branco foi, respectivamente, 16,6 e 15,8 \%.
Para todas as faixas de notas não nulas, as proporções de Concluintes foram superiores às dos Ingressantes.
As notas máximas foram 80,0, para Ingressantes, e 85,0 para Concluintes. 3.3.2.2 Comentários sobre as respostas à Questão Discursiva 38 Diversas provas obtiveram um baixo desempenho, tendo sido observado um número expressivo de respostas em branco ou destoantes em relação ao padrão de resposta esperado.
Como a questão possui dois subitens ( fatores externos e fatores internos ), nem sempre estes se apresentavam explícitos na resposta do estudante, fazendo -se necessário identificá -los durante a correção.
Outra observação a ser feita se refere às dificuldades dos estudantes quanto à elaboração de textos, já que graves erros ortográficos e de estrutura gramatical foram observados.
Destacam -se, entretanto, poucas respostas que denotavam domínio do assunto.
A presença de respostas em branco se sobrepôs às questões respondidas.
Diversas dificuldades para a elaboração de uma resposta coerente e coesa da questão 38 puderam ser identificadas.
Como por exemplo, nas respostas dos subitens, os estudantes buscaram várias tentativas para obter êxito, algumas até ousadas, mas com um resultado final distante do que foi solicitado.
Os fatores externos como os internos ( solicitados na referida questão ) foram mencionados como uma equiparação da profissão com uma atividade eminentemente humanitária, de “ amor ao próximo ”, espiritual : “ O T.O. vai para o céu ”, maternal e de caráter feminino.
Nessas situações, foi necessário ater -se ao conteúdo da resposta que, em algumas poucas provas, estava parcialmente correto.
Esforços foram empenhados pela Banca para o aproveitamento máximo das respostas apresentadas.
Por outro lado, foi possível a leitura de respostas adequadas, com conteúdo pertinente e que cumpriram corretamente o que foi pedido, essas poucas respostas, entretanto, não refletiram a grande maioria das respostas à questão de número 38. 3.3.2.3 Análise da Questão Discursiva 39 do Componente de Conhecimento Específico A Tabela 3.17 contém as informações relativas à questão 39 do conjunto de questões do componente de Conhecimento Específico.
O desempenho dos dois grupos – Ingressantes e Concluintes – nessa questão, foi superior ao desempenho na questão de número 38.
As médias foram 18,4 no grupo dos Ingressantes e 33,7 no dos Concluintes.
A proporção de estudantes que deixaram a questão em branco foi menor do que a observada na questão 38, sendo os percentuais de 44,9 \% entre os Ingressantes e de 28,0 \% entre os Concluintes.
A nota máxima ( 100,0 ) foi atingida pelos dois grupos.
Os intervalos de 30 a 40 e de 40 a 50 pontos correspondem às notas de maior freqüência ( modal ), respectivamente entre os Ingressantes com 13,2 \% e em ter os Concluintes com 15,7 \%.
As distribuições das notas entre os estudantes constam do Gráfico 3.14. 3.3.2.4 Comentários sobre as respostas à Questão Discursiva 39 Ao longo das correções, observou -se um maior número de questões respondidas e de médias melhores, comparativamente à questão de número 38.
Mesmo que parcialmente, os estudantes empenharam -se em elaborar as respostas.
Respostas corretas foram observadas, com médias, entretanto, aproximadas à nota seis.
Erros de português continuaram bastante frequentes.
Destaca -se que, em algumas provas, a ausência de caracterização do que seria pertinente a cada subitem foi notado, sendo necessário ao corretor identificar essa divisão nas respostas para evitar um prejuízo à nota dada ao estudante.
Fez -se necessário considerar ao máximo as respostas apresentadas para que fossem aproveitadas, inclusive, quando estavam objetivas, mas resumidas.
No decorrer das avaliações, houve respostas em que o estudante citou corretamente os subitens A e B ao longo da resposta, em outras, não foram mencionados ou foram subtraídos.
Destaca -se, ainda, que houve poucas trocas de questão. 3.3.2.5 Análise da Questão Discursiva 40 do Componente de Conhecimento Específico A Tabela 3.18 contém as informações relativas à questão 40 do conjunto do componente de Conhecimento Específico.
A nota média de Ingressantes foi 6,9 ; já a de Concluintes, 16,2.
A nota mediana total foi zero.
As notas máximas foram 60,0 para Ingressantes e 95,0 para Concluintes.
A questão foi deixada em branco por 60,9 \% dos Ingressantes e por 41,8 \% dos Concluintes.
O intervalo de notas não nulas no qual foi observada maior frequência, para Concluintes, corresponde à faixa que vai de 10 a 20 pontos com 13,0 \% dos estudantes.
Já para Ingressantes, houve maior incidência no intervalo entre zero e 10 pontos com 10,9 \% dos estudantes do grupo.
A distribuição de notas para os dois grupos pode ser observada no Gráfico 3.15. 3.3.2.6 Comentários sobre as respostas à Questão Discursiva 40 Quando comparada às questões anteriores, observou -se que esta questão gerou expressiva dificuldade para os estudantes.
Provavelmente, esse resultado está associado à maior complexidade das perguntas e das exigências contidas no padrão de resposta.
O número de respostas em branco ou com nota zero também foi um destaque.
Os estudantes utilizaram o artifício de mencionar trechos do texto para tentar responder, parafraseando uma tentativa de resposta.
Não conseguiram, no entanto, aferir ao trecho transposto uma análise ou comentário.
O subitem “ c ” foi o que gerou maior dificuldade para ser respondido de forma correta.
Por outro lado, o subitem “ b ” foi de o que apresentou maior índice de respostas realizadas e de acertos.
Médias superiores a seis foram raras.
Erros de português também estavam presentes mais uma vez, embora em frequência menor, porque muito do vocabulário necessário para as respostas se encontrava no próprio texto da questão.
Como ocorreu nas questões 38 e 39, em algumas provas, houve a ausência na caracterização dos subitens solicitada para responder à questão.
Nessa questão, foram expressivas as respostas parciais, somente contemplando um ou dois itens dos três solicitados para as respostas.
No subitem “ a ”, muitas respostas dos estudantes apenas retiraram trechos do texto da questão que caracterizavam as instituições totais sem comentá -los.
No padrão de resposta de parte da questão, exigia -se o comentário sobre a despersonificação ou a homegeneização, mas foi pouco apresentado nas respostas.
Alguns ainda positivaram as instituições totais.
Em relação ao subitem “ b ”, as respostas mais frequentes apontavam para a inadequação das atividades porque geravam rotina.
Muitas respostas apontavam para a vontade, o desejo ou o querer do paciente.
Também foi frequente, nas respostas dos estudantes, a apresentação, em lugar da análise da adequação ou não das atividades, com a indicação de atividades.
Pouco foi explicitado quanto à necessidade da relação das atividades com a realidade posterior, mas foi frequente relacioná -las ao seu histórico ocupacional ( “ não faziam parte do contexto e do cotidiano em que está inserido ” ).
Um número significativo de respostas também afirmava o caráter positivo das atividades para a criação de uma rotina necessária ao usuário de drogas.
Houve, ainda, respostas totalmente descontextualizadas de uma prática terapêutica, afirmando, por exemplo, a necessidade da “ fé em Deus ”.
O subitem “ c ” não apresentou nenhuma resposta praticamente em que ficasse configurada, de forma evidente, a não compulsoriedade da internação.
Muitos afirmaram que a fuga foi devido ao uso das drogas ( “ é comum a fuga, a maioria foge ” ) ou a um descontentamento com o processo de internação e das atividades realizadas.
Ainda, há respostas que apontam para o fato de a instituição não seguir as regras da ANVISA, o que teria motivado a fuga abordada no enunciado da questão.
CONSIDERAÇÕES FINAIS Na avaliação da correção, destacou -se o caráter reflexivo das questões que permitiram − através das próprias situações-problemas colocadas nos textos das questões − que o estudante criasse uma linha de raciocínio, que demonstrasse sua capacidade de responder, mesmo que parcialmente, às questões.
Considerando -se o total das três questões discursivas do componente específico, o número de questões em branco foi elevado, indicando uma participação não efetiva dos alunos na feitura da avaliação do ENADE.
O número de respostas anuladas, devido a protestos contra o ENADE, também não foi significativo, bem como as questões desconsideradas.
Foram comuns as respostas que não apresentavam conhecimentos próprios da área avaliada, sem o uso de verbetes e conceitos pertinentes.
As respostas apelavam para o senso comum e usavam termos pouco precisos.
A caracterização dos subitens que as três questões solicitavam também foi pouco realizada nas respostas.
Além disso, grande parte das respostas não desenvolveu um texto reflexivo, mas sim uma escrita “ taquigráfica ” ou em tópicos extremamente resumidos. texto reflexivo, mas sim uma escrita “ taquigráfica ” ou em tópicos extremamente resumidos.

% matched lemmas: avaliativo, crescer, excelência, grosso, interseção
\end{textsample}
